% Official source: https://www.imo-official.org/assets/documents/problems/2018/2018_eng.pdf
\competitionyear{2018}

\begin{problem}{2018}{1}
Let $\Gamma$ be the circumcircle of an acute triangle $ABC$. Points $D\in AB$ and $E\in AC$ satisfy $AD=AE$. The perpendicular bisectors of $BD$ and $CE$ meet the minor arcs $AB$ and $AC$ of $\Gamma$ at $F$ and $G$, respectively. Prove that $DE$ and $FG$ are parallel or coincide.
\end{problem}

\begin{problem}{2018}{2}
Determine all integers $n\ge3$ for which real numbers $a_1, a_2, \dots, a_{n+2}$ exist with $a_{n+1}=a_1$, $a_{n+2}=a_2$, and
\[
    a_ia_{i+1}+1=a_{i+2}\qquad(i=1, 2, \dots, n).
\]
\end{problem}

\begin{problem}{2018}{3}
An \emph{anti-Pascal triangle} is an equilateral triangular array of numbers in which every entry outside the bottom row equals the absolute difference of the two entries directly below it. For example, the following anti-Pascal triangle of four rows contains all integers from $1$ to $10$:
\[
    \begin{array}{ccccccc}
        &&&4&&&\\
        &&2&&6&&\\
        &5&&7&&1&\\
        8&&3&&10&&9
    \end{array}
\]
Does an anti-Pascal triangle with $2018$ rows exist that contains every integer from $1$ to $1+2+\cdots+2018$?
\end{problem}

\begin{problem}{2018}{4}
A \emph{site} is a point $(x, y)$ with $x, y\in\{1, 2, \dots, 20\}$. All $400$ sites start unoccupied. Amy and Ben alternate placing stones, with Amy first. Amy places a red stone on an empty site, subject to the rule that no two red stones are at distance $\sqrt5$ from one another. Ben places a blue stone on any empty site, with no distance restriction involving blue stones. The game ends as soon as a player cannot place a stone.

Find the greatest integer $K$ for which Amy can ensure she places at least $K$ red stones, regardless of Ben's moves.
\end{problem}

\begin{problem}{2018}{5}
Let $(a_n)_{n\ge1}$ be a sequence of positive integers. Suppose there exists an integer $N\ge2$ such that
\[
    \frac{a_1}{a_2}+\frac{a_2}{a_3}+\cdots+\frac{a_{n-1}}{a_n}+\frac{a_n}{a_1}
\]
is an integer for all $n\ge N$. Prove that some positive integer $M$ satisfies $a_m=a_{m+1}$ for every $m\ge M$.
\end{problem}

\begin{problem}{2018}{6}
A convex quadrilateral $ABCD$ satisfies $AB\cdot CD=BC\cdot DA$. A point $X$ inside $ABCD$ satisfies
\[
    \angle XAB=\angle XCD, \qquad \angle XBC=\angle XDA.
\]
Prove that $\angle BXA+\angle DXC=180^\circ$.
\end{problem}
